\color{unidarkgreen}\section[Использование по назначению]{Использование по назначению}
\color{black}

\needspace{4\baselineskip}
\color{unidarkgreen}{\subsection{Эксплуатационные ограничения}}\color{black}

\begin{enumerate}[label=\arabic{section}.\arabic{subsection}.\arabic*, labelsep=4pt, leftmargin=0pt, itemindent=57pt]
    
\item
Эксплуатация и обслуживание устройства должны проводиться в соответствии с НТД РФ - «Правилами технической эксплуатации электрических станций и сетей РФ», «Правил по охране труда при эксплуатации электроустановок», «Правил устройств электроустановок» и требованиями настоящего РЭ.

\item
Возможность работы устройства в условиях, отличных от указанных в РЭ, должна согласовываться с предприятием-держателем подлинников конструкторской документации и с предприятием-изготовителем.

\item
Условия эксплуатации в части воздействия механических факторов должны соответствовать требованиям настоящим РЭ.

\end{enumerate}

\color{unidarkgreen}{\subsection{Подготовка к эксплуатации}}\color{black}
\subsubsection{Меры безопасности}

\begin{enumerate}[label=\arabic{section}.\arabic{subsection}.\arabic{subsubsection}.\arabic*,  align=left, labelsep=4pt, leftmargin=0pt, itemindent=57pt]

\item
Монтаж, обслуживание и эксплуатацию шкафа разрешается производить лицам, прошедшим специальную подготовку, имеющим аттестацию на право выполнения работ (с учётом соблюдения необходимых мер защиты изделий от воздействия статического электричества), хорошо знающим особенности электрической схемы и конструкцию шкафа.

\item
Монтаж шкафа и работы на разъёмах терминала, рядах зажимов шкафа и разъёмах устройств следует производить при обесточенном состоянии шкафа.
При необходимости проведения проверок с наличием напряжения, должны приниматься дополнительные меры, предотвращающие поражения обслуживающего персонала электрическим током.

\item
По требованиям защиты человека от поражения электрическим током шкаф соответствует классу I по ГОСТ 12.2.007.0-75.

\item
Запрещается подача какого-либо электроснабжения без заземления.
Шкаф перед подачей электроснабжения и во время работы должен быть надежно заземлен.

\end{enumerate}

%\subsubsection{Установка и монтаж ВЧ-приемопередатчика} %для шкафов с функциями ДФЗ/НВЧЗ/ВЧБ
%
%\begin{enumerate}[label=\arabic{section}.\arabic{subsection}.\arabic{subsubsection}.\arabic*,  align=left, labelsep=4pt, leftmargin=0pt, itemindent=57pt]
%
%\item
%ВЧ-пост перед установкой должна быть проверена и отрегулирована в соответствии с заводской инструкцией.
%
%\item
%Монтаж и подключение ВЧ-поста к шкафу проводится непосредственно на месте эксплуатации шкафа в соответствии со схемой и заводской инструкцией.
%Установка ВЧ-поста производится в предусмотренное заводом место, согласно чертежа общего вида (ВО).
%
%\item
%Подключение цепей связи релейной части защиты с приёмопередатчиком осуществляется проводами, связанными в жгут.
%На проводах имеется маркировка.
%Необходимо развязать жгут и соединить провода с выходными зажимами приёмопередатчика.
%
%\end{enumerate}

\subsubsection{Монтаж шкафа}

\begin{enumerate}[label=\arabic{section}.\arabic{subsection}.\arabic{subsubsection}.\arabic*,  align=left, labelsep=4pt, leftmargin=0pt, itemindent=57pt]

\item
Шкаф не подвергается консервации смазками и маслами и расконсервации не подлежит.

\item
Выполнить подключение шкафа согласно утвержденному проекту в соответствии с указаниями настоящего РЭ. 

\item
Связь шкафа с другими шкафами защит и устройствами производить с помощью кабелей или проводников в соответствие с принятыми проектными техническими решениями.

\item
Ввод внешних кабелей осуществляется через гермовводы, снизу.
С целью устранения механического усилия тяжения кабеля в шкафу (требование п.2.1.24 ПУЭ), входящие в шкаф кабели вторичных цепей на входе должны быть закреплены зажимом кабельным (к устройству крепления и заземления экранов кабелей) или вводом кабельным с контргайкой типа PG/MG (к панели ввода).

\item
Монтаж заземления экранов внешних кабелей необходимо проводить после установки и закрепления шкафа на конструкциях, предусмотренных проектом, и прокладки всех контрольных кабелей.

\item
Для заземления экрана кабеля рекомендуется использовать кабельные хомуты из нержавеющей стали или экранирующие зажимы.
Кабельные хомуты должны максимально охватывать всю наружную электропроводящую поверхность экрана кабеля и соответствующую этому кабелю перемычку устройства заземления, обеспечивая между ними надёжный электрический контакт.

\item
Гермовводы на днище шкафа предназначены для ввода кабелей максимальным диаметром до 20 мм.
В случае иных требованиях, данные требования необходимо указать в картах заказа.

\item
Подключение высокочастотной части защиты к высокочастотному каналу связи производить с помощью коаксиального кабеля непосредственно к выходным зажимам приёмопередатчика.

\end{enumerate}

\subsubsection{Подготовка шкафа к работе}

\begin{enumerate}[label=\arabic{section}.\arabic{subsection}.\arabic{subsubsection}.\arabic*,  align=left, labelsep=4pt, leftmargin=0pt, itemindent=57pt]

\item
Шкаф выпускается с предприятия работоспособным и полностью испытанным.
Проверка монтажа и целостности электрических цепей шкафа проведена на предприятии изготовителе.

\item
Положение оперативных переключателей и значения уставок защит шкафа выставить с учётом бланка уставок.

\item
Настройки терминала и необходимая оперативная информация доступны на дисплее терминалы, при входе в соответствующие пункты меню, с помощью нажатия на соответствующие кнопки управления.
Работа с терминалом подробно описана в руководстве на установленный терминал.

\end{enumerate}

\color{unidarkgreen}\subsection{Использование изделия}\color{black}

\begin{enumerate}[label=\arabic{section}.\arabic{subsection}.\arabic*,  align=left, labelsep=4pt, leftmargin=0pt, itemindent=57pt]

\item
При вводе шкафа в эксплуатацию необходимо выполнить следующие работы:
\begin{itemize}
\item проверку переходного сопротивления заземления шкафа;
\item проверку сопротивления и прочности изоляции;
\item проверку коммутационного оборудования в составе шкафа;
\item задание уставок и конфигурации защит терминала в составе шкафа;
\item проверку правильности отображения аналоговых величин, поданных от постороннего источника;
\item проверку уставок (ИО) защит устройства;
\item проверку воздействия на внешние цепи;
\item проверку действия на центральную сигнализацию;
\item проверку взаимодействия шкафа с другими устройствами;
\item проверку шкафа рабочим током и напряжением.
\end{itemize}

\item
По результатам проверок составляется протокол(ы) испытания шкафа при вводе в эксплуатацию (протокол проверок при новом включении).

\end{enumerate}

\color{unidarkgreen}\subsection{Возможные неисправности и методы их устранения}\color{black}

\begin{enumerate}[label=\arabic{section}.\arabic{subsection}.\arabic*,  align=left, labelsep=4pt, leftmargin=0pt, itemindent=57pt]

\item
Неисправности могут возникнуть при нарушении условий транспортирования, хранения и эксплуатации.

\item
При включении питания и в процессе работы шкафа могут возникнуть неисправности, обнаруживаемые системой контроля терминала.
Описание возможных неисправностей терминала и методов их устранения приведены в РЭ на терминал.

\item
Обслуживающий шкаф персонал может самостоятельно провести ремонт или замену коммутационных устройств, светосигнальной арматуры и т.д.

% ПРЕДУПРЕЖДЕНИЕ
%\needspace{4\baselineskip}
\medskip
\noindent
\setlength{\fboxrule}{1.2pt}
\setlength{\fboxsep}{3pt}
\noindent
\fcolorbox{unidarkgreen}{white}{%
  \begin{minipage}{\dimexpr\linewidth-2\fboxrule-2\fboxsep}
    {\hspace*{0.5em}\textcolor{red!80!black}{\Large\faExclamationTriangle}}\hspace{0.5em}
    \hyphenpenalty=10000  \textbf{В СЛУЧАЕ ОБНАРУЖЕНИЯ НЕИСПРАВНОСТЕЙ В ШКАФУ, НЕОБХОДИМО НЕМЕДЛЕННО ПОСТАВИТЬ В ИЗВЕСТНОСТЬ ПРЕДПРИЯТИЕ-ИЗГОТОВИТЕЛЬ.}
  \end{minipage}%
}\vspace{1mm}

%\item
%\textbf{В СЛУЧАЕ ОБНАРУЖЕНИЯ НЕИСПРАВНОСТЕЙ В ШКАФУ, НЕОБХОДИМО НЕМЕДЛЕННО ПОСТАВИТЬ В ИЗВЕСТНОСТЬ ПРЕДПРИЯТИЕ-ИЗГОТОВИТЕЛЬ.}

\end{enumerate}